\scanpage{80}\mspage{67}
% archive sheet number: 79

Mais il \unsure{se|pr} présente à nouveau la perplexité
si le foncteur (3.3.1) ainsi construit
est bien un adjoint à gauche de $u_W^*$.
On peut songer à utiliser encore la
$\unsure{=}$ factorisation en $u=vw$, en utilisant
le mapping-cône. La difficulté, cependant,
c'est qu'à présent $J=C(u)$ n'est plus une
cat ordonnée, donc on n'a pas une
définition toute prête de \deleted{\illegible{1}} $C_{J^{\unsure{o|\circ}}}$
$\subset\mathrm{Fl}\, M(J^o)$, donnant lieu à
une sous-catégorie $M(\deleted{\unsure{\in}}J^o)_c$ etc. En
fait, il semblerait qu'on soit pratiquement
paralysé, \unsure{tant} qu'on n'arrive à construire
sur $M(J^o)$ également une structure
de catégorie à \unsure{cofibration|cofibrations}, se comportant
bien par rapport au foncteur $\unsure{v|w}_!$.
La ``démonstration'' d'Anderson de l'existence
de $u_!$ pour \uline{\unsure{tt}} flèches dans $\mathrm{Cat}$,
foire totalement, \deleted{\unsure{car}} ($=$ pour $I\unsure{\text{---}|\to}e$,
$I$ un ordonné fini --- mais là j'ai
pu rétablir la situation), car il se ramène
au cas où $I$ est ordonné loc. fini, et il
